\documentclass[11pt]{article}
\title{How the Codynamic Book Machine Works}
\author{Arthur Petron}
\usepackage[margin=1in]{geometry}
\usepackage{graphicx}
\usepackage{amsmath,amssymb,mathtools}
\usepackage{tikz}
\usetikzlibrary{positioning,arrows.meta}
\usepackage{hyperref}
\graphicspath{{media/}{media/diagrams/}{images/}{artwork/}}

\begin{document}

\section*{Schema, Intent, and Metadata}

The Codynamic Book Machine begins by capturing the \emph{schema} of a work before it generates any prose. That choice is not merely administrative. In a multi-agent writing system, the earliest representation of a book must already answer the questions that later agents will need in order to act consistently: What is this work for? Who is it for? What voice should it use? What publication context will it live in? Which artifacts belong to it? Without those answers, generation becomes guesswork, and guesswork is expensive once drafts, citations, diagrams, and builds begin to accumulate.

The canonical work object therefore stores intent alongside structure. Intent fields describe the purpose of the project, the intended audience, the reader takeaway, the authorial voice, and the genre or writing style. These fields are not decorative metadata; they are operational constraints. A section drafted for researchers should sound different from one drafted for general readers. A system described as a technical exposition should emphasize mechanisms, data flow, and design tradeoffs rather than narrative flourish. By recording intent up front, the machine gives downstream agents a stable target for tone, scope, and level of detail.

In practice, this means the work record can carry values such as \texttt{purpose}, \texttt{audience}, \texttt{reader\_takeaway}, \texttt{author\_persona}, \texttt{writing\_style}, and \texttt{genre}. Those fields let the system distinguish between what the book is trying to accomplish and how it should speak while doing so. A draft can then be checked against the declared brief instead of against an informal memory of the brief. That distinction matters because the same outline can support very different books depending on whether the intended reader is a researcher, a tool builder, or a general technical audience.

Metadata serves a similar function at the publication layer. Version, creation and update timestamps, language, license, maintenance responsibility, and status all help the system distinguish a working draft from a publishable artifact. Publication settings such as document class, paper size, and front-matter options also belong in the canonical record because they affect compilation and output behavior. In other words, the work object is not just a content container; it is the source of truth for how the content should be interpreted, edited, and rendered.

A simple intake flow makes this concrete. Before a book is decomposed into chapters and sections, the system asks for the title, purpose, audience, reader takeaway, voice, genre, and design expectations, then normalizes those answers into the canonical work object. Once that happens, later agents no longer have to infer intent from prose alone. They can read it directly from the schema and act on it consistently. The same principle applies to publication metadata: a draft that is still under active revision should not be treated as a finished release, and a work intended for a book-class PDF should not be compiled as if it were a slide deck or article.

Capturing this information before generation also improves coordination. Section agents can read the same intent fields and produce prose that remains aligned across chapters. Gardener and design agents can evaluate whether a draft matches the declared audience and style. Reference and diagram workflows can be triggered only when the section context indicates that they are needed. The result is a system in which authorial intent is explicit, machine-readable, and available to every participant in the workflow.

This schema-first approach also reduces drift. If purpose, audience, and publication metadata are embedded in the work object, then later revisions do not have to reconstruct them from scattered notes or informal memory. The machine can validate the work against its declared shape, preserve continuity across agent passes, and compile outputs that remain faithful to the original brief. In that sense, schema is not a preliminary formality; it is the contract that makes codynamic authoring possible.

Publication metadata matters for the same reason. By storing status, language, license, and output settings in the work record, the system keeps editorial intent and rendering intent in the same place. That makes the work object a practical control surface rather than a passive archive. It also gives the runtime a reliable way to distinguish a draft from a release candidate, and a book from some other output form, before any expensive generation or compilation work begins.

In short, schema comes first because every later step depends on it. Intent tells the machine what to say, metadata tells it how to package and maintain the result, and the canonical work object gives every agent a shared source of truth. The rest of the system---outline decomposition, section drafting, review, diagrams, references, and compilation---can only remain coherent if they all begin from that same explicit foundation.


\end{document}
